%% kapa: exam=1 round=23 year=2012 period=1 form=A subject=economics subject_ko=경제원론 seq=16 number=56
%% source: Q-Net 감정평가사 기출문제 (cst00309) artlSeq=5052578 file=raw/1008/5052578/2012년 제23회 감정평가사 제1차시험 기출문제(1교시 A형).pdf
%% conversion: kordoc 4.24.51 -> tools/gen1.py (auto); review=proofread-sonnet; flags=none
\begin{question}{56}
외부성에 관한 설명으로 옳은 것을 모두 고른 것은?
\begin{framedtext}
ㄱ. 부(-)의 외부성이 존재하면 시장 생산량은 사회적 최적 생산량보다 많다. \par
ㄴ. 외부성은 합병이나 보조금 혹은 조세 등을 통해 내부화시킬 수 있다. \par
ㄷ. 코우즈(R. Coase)에 의하면 외부성이 존재하더라도 재산권이 확립되어 있으면 정부의 개입이 불필요할 수 있다.
\end{framedtext}
\begin{choices}
  \item ㄱ
  \item ㄱ, ㄴ
  \item ㄴ, ㄷ
  \item ㄱ, ㄷ
  \item ㄱ, ㄴ, ㄷ
\end{choices}
\end{question}
\answer{5}
